\clearpage
\chapter*{Errata da exposição n\textsuperscript{o} 182}
\addcontentsline{toc}{chapter}{Errata da exposição n\textsuperscript{o} 182}
\setcounter{footnote}{1}
\src{298}

\noindent\textbf{2. Geometria formal e geometria algébrica.}

\begin{center}
[Seminário Bourbaki, t. 11, 1958/59, n\textsuperscript{o} 182, 28 p.]
\end{center}

\noindent\textbf{Página 8, teorema 6.} --- A hipótese «$f$
quase-projetivo» pode ser substituída pela hipótese mais fraca «$f$
separado», em virtude do resultado seguinte (cf. SGA VIII, 6.2): todo morfismo
$f:X\longrightarrow Y$ que seja quase-finito e separado é quase-projetivo.

\src{299}
\noindent\textbf{Página 12, e página 14, observação 1.} --- Assinalemos que J.-P.
Serre construiu uma variedade projetiva não singular, de dimensão 3, sobre
um corpo algebricamente fechado $k$, de característica $p>0$, que não provém
por redução de um esquema próprio sobre um anel local que seja domínio de
integridade, com corpo residual $k$ e corpo de frações de característica
zero.\footnote{Serre (Jean-Pierre). --- Exemples de variétés projectives en
caractéristique $p$ non relevables en caractéristique zéro, \emph{Proc. Nat.
Acad. Sc. U.S.A.}, t. 47, 1961, p. 108--109.} Mumford teria encontrado um exemplo
análogo, com uma superfície projetiva não singular.

\noindent\textbf{Página 15, observações 2 e 3.} --- Atualmente sou menos
otimista quanto aos resultados aqui conjecturados. No entanto, o
problema relativo às variações da estrutura do espaço projetivo,
assinalado ao final da observação 3, foi resolvido afirmativamente por
Hironaka, e o problema análogo para as variedades abelianas foi resolvido
por Koizumi.

\noindent\textbf{Página 16, linha 3.} --- Assinalemos que Mumford construiu
recentemente os esquemas de módulos para as curvas de gênero $g$ (cf. SHMI).
Além disso, o teorema 10 prova que os «esquemas de Jacobi de nível $n$»
da teoria dos módulos são não singulares (e até mesmo lisos sobre $\mathbf Z$).

\noindent\textbf{Página 23, linha 12, primeira linha do corolário
5.} --- Acrescentar: «sendo $X$ ou $Y$ próprio sobre $k$».

A substância dos §§ 1 a 5 está contida na parte publicada de EGA III, e
a dos §§ 6 e 7 está contida em SGA III. Para o estudo do grupo
fundamental, vejam-se SGA V, IX, X, XI, bem como SGA 1962 (exposições X, XII e
XIII), para os teoremas do tipo de Lefschetz e numerosas questões em aberto.
Somente a teoria dos recobrimentos moderadamente ramificados (cf. teorema 14)
ainda não foi objeto de uma redação formal. O corolário do teorema 14,
que fornece a determinação completa dos recobrimentos galoisianos de ordem
coprima à característica de uma curva algébrica sobre um corpo
algebricamente fechado, foi utilizado de maneira essencial em três ocasiões:

1\textsuperscript{o} na demonstração, por Igusa, da desigualdade de Picard
para as superfícies projetivas não singulares em característica arbitrária;

2\textsuperscript{o} no estudo (desenvolvido independentemente por Ogg e
Šafarevič) do grupo dos fibrados principais homogêneos sob uma variedade
abeliana definida sobre um corpo de funções de uma variável, em
característica arbitrária;

3\textsuperscript{o} na demonstração recente, por Artin, de certos
teoremas fundamentais sobre a «cohomologia de Weil» das variedades algébricas.
